\section{Conclusion}
LoRA offers a practical way to reduce the number of parameters updated
during fine-tuning. For a $1024 \times 1024$ weight matrix, rank-4 LoRA
uses 8,192 trainable parameters instead of 1,048,576, a 99.22\%
reduction. In the DistilRoBERTa experiment on AG News, the
trainable adapters and classifier account for 0.8944\% of the model's
parameters, and the model reaches 93.76\% test accuracy and 93.76\%
macro-F1 across 7,600 examples. The full fine-tuning baseline reaches
95.03\% accuracy and macro-F1, about 1.26 percentage points higher on
both measures. LoRA uses 47.31 rather than 66.38 minutes training wall
time and peaks at 1.099 rather than 2.138 GiB allocated GPU memory.
These runs show a measured quality and resource trade-off.

The comparison has only one run per method, different learning rates,
and no comparison across ranks; the LoRA result files do not identify
the rank.
Freezing the base weights reduces gradient and optimizer-state
requirements, but the weights, activations, and intermediate
computations still consume memory. Rank selection also trades
adaptation capacity against parameter cost, while keeping adapters
separate or merging them affects deployment flexibility. O-LoRA
extends the adapter approach by encouraging task-update subspaces to
be orthogonal in continual learning, but it was not evaluated in our
experiment. Overall, LoRA is most appropriate when limiting trainable
parameters is important; its quality, resource benefits, and
suitability should be verified under the target task and deployment
conditions.
